\documentclass[10pt]{article}

\usepackage{amsmath,amssymb,amsthm}

\newtheorem{definition}{Definition}
\newtheorem{problem}{Problem}
\newtheorem{theorem}{Theorem}
\newtheorem{remark}[theorem]{Remark}

\newcommand{\SRI}{\text{SR-I}}
\newcommand{\SRII}{\text{SR-II}}
\newcommand{\SRIII}{\text{SR-III}}
\newcommand{\SRCI}{\text{SRC-I}}
\newcommand{\SRCII}{\text{SRC-II}}
\newcommand{\SRCIII}{\text{SRC-III}}
\newcommand{\Res}{\text{Res}}

\hyphenation{Boo-le-an}

\title{Eliminating Weakening in Resolution with Dynamic Symmetries - Hard}
\author{Nikita Gaevoy}
\date{}

\begin{document}

\maketitle

\section{Preliminaries}

We work with propositional formulas in conjunctive normal form over the
Boolean variables $x_1,\ldots,x_n$. A \emph{literal} is a variable $x$ or a
negated variable $\lnot x$, a \emph{clause} is a finite disjunction of
literals, identified with the set of the literals it contains, and a
\emph{CNF formula} is a finite conjunction of clauses, identified with the
set of the clauses it contains. We write
\[
\operatorname{size}(\Phi)
\]
for the total number of literal occurrences in $\Phi$. The empty clause is
denoted by $\bot$.

The \emph{resolution rule} is
\[
\frac{
A\lor x
\qquad
B\lor\lnot x
}{
A\lor B
},
\]
and the variable $x$ is called the \emph{pivot} of the inference. A
\emph{Resolution refutation} of an unsatisfiable CNF formula $\Phi$ is a
sequence of clauses in which every clause is either a clause of $\Phi$ or is
obtained from two earlier clauses by the resolution rule, and whose last
clause is $\bot$.

\subsection{Small symmetries}

\begin{definition}[$k$-symmetry]\label{def:ksym}
Let $k\ge 1$ be either a constant or a function of the number of variables
$n$. A mapping $\varphi$ of literals into literals is a \emph{$k$-symmetry} if
and only if it satisfies the following two conditions:
\begin{enumerate}
    \item $\varphi(\lnot x)=\lnot\varphi(x)$ for every variable $x$;
    \item all variables except for at most $k$ are mapped into themselves,
    i.e.\ $\varphi(x)=x$ for all but at most $k$ variables $x$.
\end{enumerate}
A $k$-symmetry is \emph{positive} if it maps variables into variables without
negation. Symmetries that are not necessarily positive are called
\emph{complementary} symmetries.

A $k$-symmetry is extended to clauses and to sets of clauses pointwise:
\[
\varphi(C):=\{\varphi(l):l\in C\},
\qquad
\varphi(\Phi):=\{\varphi(C):C\in\Phi\}.
\]
\end{definition}

\begin{remark}
As usual for symmetries, $\varphi$ is required to be a bijection of the set of
literals; by condition~2 of Definition~\ref{def:ksym} it is therefore given by
a permutation of at most $k$ variables together with a choice of a sign for
each of them. Dropping the bijectivity requirement yields the notion of a
$k$-homomorphism and the corresponding proof systems
$\text{HR-I}$, $\text{HR-II}$, $\text{HR-III}$ of~\cite{Szeider05}, which we
do not consider here.
\end{remark}

When $k$ is equal to the number of variables, condition~2 is vacuous and one
recovers the classical notion of a symmetry of~\cite{Krishnamurthy85,Urquhart99};
in this case the parameter $k$ is omitted from the notation.

\begin{definition}[Resolution with dynamic $k$-symmetries]\label{def:srciii}
Let $k\ge 1$ be a constant or a function of the number of variables.
\emph{Resolution with dynamic $k$-symmetries}, denoted $k$-$\SRCIII$, is the
proof system with the following two rules.
\begin{enumerate}
    \item \textbf{Resolution rule.}
    \[
    \frac{
    A\lor x
    \qquad
    B\lor\lnot x
    }{
    A\lor B
    }.
    \]
    \item \textbf{Symmetry rule.}
    \[
    \frac{
    \varphi(B_1)
    \qquad
    \varphi(B_2)
    \qquad
    \cdots
    \qquad
    \varphi(B_l)
    }{
    \varphi(A)
    }\ \varphi,
    \]
    where $\varphi$ is a $k$-symmetry and $A$ was inferred from
    $B_1,B_2,\ldots,B_l$ earlier in the proof. Each time this rule is applied,
    the symmetry $\varphi$ is written down as a part of the proof, for the
    purposes of verification. We call the inference of $A$ from
    $B_1,\ldots,B_l$ the \emph{justification} of the symmetry rule and
    $\varphi(A)$ the \emph{resulting clause}.
\end{enumerate}
As in Resolution, a proof is a sequence of clauses in which every clause is
either an input clause or is inferred from previous clauses by one of the two
rules above, and the last clause is $\bot$.
\end{definition}

The motivation behind the symmetry rule is the following. If $A$ was inferred
from $B_1,\ldots,B_l$, and $\varphi$ is a $k$-symmetry mapping all the premises
to clauses that have already been derived, then applying $\varphi$ to that
inference yields a valid derivation of $\varphi(A)$ from
$\varphi(B_1),\ldots,\varphi(B_l)$. Hence, instead of repeating the derivation, one
may infer $\varphi(A)$ in a single step.

The system of Definition~\ref{def:srciii} is \emph{dynamic}: the premises
$B_i$ of the symmetry rule are arbitrary clauses of the proof and need not
belong to the input formula. Restricting them to the input formula gives the
\emph{static} systems.

\begin{definition}[Static systems]
\emph{Resolution with local $k$-symmetries} ($k$-$\SRCII$) is
$k$-$\SRCIII$ with the additional requirement that in every application of
the symmetry rule all $B_i$ and all $\varphi(B_i)$ are clauses of the input
formula.

\emph{Resolution with global $k$-symmetries} ($k$-$\SRCI$) is
$k$-$\SRCIII$ with the additional requirement that in every application of
the symmetry rule both sets $\{B_i\}$ and $\{\varphi(B_i)\}$ are equal to the set
of all input clauses.

\emph{$k$-$\SRI$, $k$-$\SRII$, and $k$-$\SRIII$} are the systems
$k$-$\SRCI$, $k$-$\SRCII$, and $k$-$\SRCIII$, respectively, with the
additional requirement that all symmetries used are positive.
\end{definition}

If $n$ is the number of variables, then $n$-$\SRCIII$, $n$-$\SRIII$ are the
systems $\SRCIII$, $\SRIII$ of Szeider~\cite{Szeider05}, while $n$-$\SRCI$,
$n$-$\SRCII$ are the systems $\SRCI$, $\SRCII$ of Urquhart~\cite{Urquhart99}
and $n$-$\SRI$, $n$-$\SRII$ are the systems of
Krishnamurthy~\cite{Krishnamurthy85} and of Arai and Urquhart~\cite{AU00}.

All of the systems above are Cook--Reckhow proof systems~\cite{CR79}:
completeness follows from the fact that they extend Resolution, polynomial-time
verification is possible because the symmetry used is written down at every
application of the symmetry rule, and soundness follows from the requirement
that the justification be derived before the symmetry rule is applied.

\begin{definition}[Proof size]
For a proof system $\Pi$ and an unsatisfiable CNF formula $\Phi$, we write
\[
S_{\Pi}(\Phi)
\]
for the number of steps in the shortest $\Pi$-proof of $\Phi$. (For all the
systems considered here, the number of steps and the total bit length of a
proof differ by at most a polynomial factor.)
\end{definition}

Note that, because of the symmetry rule, $S_{\Pi}(\Phi)$ may be smaller than
the number of clauses of $\Phi$.

Following~\cite{Gaevoy26}, the definition of $k$-$\SRCIII$ above does not
include the weakening rule
\[
\frac{A}{A\lor x}.
\]
We write $k$-$\SRCIII^{w}$ for the system obtained from $k$-$\SRCIII$ by
adding this rule. For the static systems the analogous distinction is
immaterial: weakening can be eliminated exactly as in
Resolution~\cite{Szeider05}, so for all $k$ the systems $k$-$\SRI$,
$k$-$\SRCI$, $k$-$\SRII$, and $k$-$\SRCII$ are unaffected by its addition. For
the dynamic systems no such elimination is known, and the addition of
weakening may potentially make the system stronger.

\begin{definition}[p-simulation]\label{def:psim}
A proof system $\Pi_1$ \emph{p-simulates} a proof system $\Pi_2$ if there is a
polynomial $q$ such that
\[
S_{\Pi_1}(\Phi)
\ \le\
q\bigl(S_{\Pi_2}(\Phi)+\operatorname{size}(\Phi)\bigr)
\]
for every unsatisfiable CNF formula $\Phi$, and the two systems are
\emph{p-equivalent} if each p-simulates the other. Since every
$k$-$\SRCIII$-proof is a $k$-$\SRCIII^{w}$-proof, the system
$k$-$\SRCIII^{w}$ p-simulates $k$-$\SRCIII$ for trivial reasons; the content
of the notion here is the converse direction, that is, the elimination of the
weakening rule.
\end{definition}

\section{Problem formulation}

\begin{problem}[Elimination]\label{prob:elim}
Prove that $k$-$\SRCIII$ p-simulates $k$-$\SRCIII^{w}$, i.e.\ that every
application of the weakening rule can be eliminated with at most a polynomial
increase in proof size, as is the case for Resolution and for all the static
systems $k$-$\SRI$, $k$-$\SRCI$, $k$-$\SRII$, $k$-$\SRCII$~\cite{Szeider05}.
\end{problem}

\begin{remark}\label{rem:dichotomy}
The complementary statement --- that the two systems are separated, i.e.\ that
there is a family $(\Psi_n)_{n\ge 1}$ with
$S_{k\text{-}\SRCIII^{w}}(\Psi_n)=\operatorname{poly}(\operatorname{size}(\Psi_n))$
and $S_{k\text{-}\SRCIII}(\Psi_n)=\operatorname{size}(\Psi_n)^{\omega(1)}$
--- is a separate problem. Exactly one of the two
is true for a given $k$, and both are open: for every $k\ge 1$, including
$k=1$, and also for the unrestricted system $\SRCIII$, that is, for $k$ equal
to the number of variables.
\end{remark}

\begin{remark}\label{rem:independent}
Problem~\ref{prob:elim} asks for a transformation of proofs, not for a hard
formula, and a solution to it neither uses nor implies any lower bound. The
separating family of Remark~\ref{rem:dichotomy} is the opposite case: it
witnesses a superpolynomial lower bound for $k$-$\SRCIII$, and no such bound is
known for any $k$.

The implication runs the useful way. A solution to Problem~\ref{prob:elim}
makes the two readings of a proof-size lower bound for $k$-$\SRCIII$ coincide:
a bound proved without the weakening rule would then transfer to
$k$-$\SRCIII^{w}$ for free, so a solver could assume without loss of
generality that no weakening occurs.
\end{remark}

\begin{remark}\label{rem:scope}
The scope is the dynamic system only. For the static systems the question does
not arise: weakening can be eliminated there exactly as in Resolution, so for
all $k$ the systems $k$-$\SRI$, $k$-$\SRCI$, $k$-$\SRII$ and $k$-$\SRCII$ are
unaffected by its addition~\cite{Szeider05}. What breaks in the dynamic case is
that the premises of the symmetry rule are arbitrary clauses of the proof, so
a weakening applied to a justification cannot simply be pushed to the leaves.
\end{remark}

\begin{thebibliography}{Gae26}

\bibitem[AU00]{AU00}
Noriko~H. Arai and Alasdair Urquhart.
\newblock Local symmetries in propositional logic.
\newblock In {\em Automated Reasoning with Analytic Tableaux and Related
  Methods (TABLEAUX 2000)}, pages 40--51. Springer, 2000.

\bibitem[CR79]{CR79}
Stephen~A. Cook and Robert~A. Reckhow.
\newblock The relative efficiency of propositional proof systems.
\newblock {\em Journal of Symbolic Logic}, 44(1):36--50, 1979.

\bibitem[Gae26]{Gaevoy26}
Nikita Gaevoy.
\newblock The power of small symmetries.
\newblock In {\em 51st International Symposium on Mathematical Foundations of
  Computer Science (MFCS 2026)}, LIPIcs. Schloss Dagstuhl -- Leibniz-Zentrum
  f\"{u}r Informatik, 2026.
\newblock To appear.

\bibitem[Kri85]{Krishnamurthy85}
Balakrishnan Krishnamurthy.
\newblock Short proofs for tricky formulas.
\newblock {\em Acta Informatica}, 22(3):253--275, 1985.

\bibitem[Sze05]{Szeider05}
Stefan Szeider.
\newblock The complexity of resolution with generalized symmetry rules.
\newblock {\em Theory of Computing Systems}, 38(2):171--188, 2005.

\bibitem[Urq99]{Urquhart99}
Alasdair Urquhart.
\newblock The symmetry rule in propositional logic.
\newblock {\em Discrete Applied Mathematics}, 96--97:177--193, 1999.

\end{thebibliography}

\end{document}
